\section{What a BSP Should Expose}
\label{sec:bsp-api}

A BSP's public API is a vocabulary, and the test of a vocabulary is what a user can say with it without
knowing anything else.

\subsection{The Right Shape}

\begin{minted}{rust}
let module = Module::take().expect("peripherals taken twice");
let mut led = module.led();

let base = ExpansionBase::take().expect("peripherals taken twice");
let mut oled = base.oled();
let env = base.environment()?;      // BME68x
let rtc = base.rtc()?;              // PCF8563
let mut buzzer = base.buzzer();
\end{minted}

Every one of those names is a thing on the desk. None of them is an address, a bus, a pin, an alias or
a register — and that is the whole specification.

\subsection{What the Signature Hides, and What It Must Not}

\begin{table}[htbp]
    \centering
    \footnotesize
    \begin{tabular}{@{}p{0.26\textwidth}p{0.30\textwidth}p{0.40\textwidth}@{}}
        \toprule
        \textbf{Hidden} & \textbf{Example} & \textbf{Why hiding it is right} \\
        \midrule
        addresses & \codeinline{0x3C}, \codeinline{0x76}, \codeinline{0x51} & a user cannot check them and must not need to \\
        pins & \codeinline{PC04}/\codeinline{PC05}, \codeinline{PB2}/\codeinline{PB3} & measured facts; wrong in half the documents \\
        power gates & \codeinline{PD5}, \codeinline{PC8} & sequencing the user would get wrong \\
        aliases & which peripheral window answered & a security-attribution detail \\
        baud and clock & I2C half-period, PWM frequency & driver knowledge \\
        \bottomrule
    \end{tabular}
    \caption{What a facade should hide. Note that every row is a \emph{fact}, not a \emph{choice} — the
    user has no decision to make about any of them.}
    \label{tab:bsp-hides}
\end{table}

\subsection{What It Must Not Hide: Structure}

The temptation is to hide the two buses, because ``two I2C buses'' is an implementation detail. It is
not — it is a fact about the hardware that changes what the user can do:

\begin{itemize}
    \item The panel, the BME68x and the RTC share \codeinline{PC04}/\codeinline{PC05}.
    \item The IMU has its own pair, \codeinline{PB2}/\codeinline{PB3}, because the module's sensors are
          wired to the module's pins.
    \item A single \codeinline{i2c()} would have to take a ``which bus'' argument anyway, so hiding it
          only moves the choice to a place with less information.
\end{itemize}

The rule that came out of this: \textbf{hide the addresses, keep the structure}. A facade that hides
structure has to guess, and guessing is what the layer below exists to prevent.

\subsection{Conflicts Should Be Refused, Not Resolved}

\begin{minted}{rust}
// The microSD card's chip select shares a pin on this board. When that driver lands,
// this facade is where the conflict is refused -- not silently arbitrated.
pub fn sd(&mut self) -> Result<SdCard<'_>, BoardError>
\end{minted}

A facade that resolves a real conflict silently is a facade that produces a bug the user cannot see.
Returning an error — or not offering the operation at all — is honest, and it is the difference between
an API that documents the hardware and one that pretends the hardware is simpler than it is.
